% ===========================================================================
%  CAPÍTULO 11 — RESULTADOS
%  Trabajo de Fin de Grado · Ingeniería del Software
% ===========================================================================

\section*{11. Resultados}
\label{sec:resultados}

Como se detalló en el apartado 7.3.3, para evaluar la calidad
predictiva de los modelos desarrollados, se establecieron umbrales
de rendimiento basados en la distribución del Error Absoluto Medio
($\mathrm{MAE}$) observada en el estudio previo. Estos objetivos
permiten clasificar la precisión de los modelos según la posición
táctica.

% ---------------------------------------------------------------------------
\subsubsection{Objetivos de rendimiento por posición}
\label{subsubsec:objetivos_rendimiento}

\begin{table}[H]
\centering
\small
\begin{tabular}{|l|c|c|c|}
\hline
\textbf{Posición} & \textbf{Élite} & \textbf{Aceptable} &
\textbf{Alerta} \\
\hline
Porteros    (PT) & $\le 2{,}93$ & $\le 3{,}18$ & $> 3{,}63$ \\
Defensas    (DF) & $\le 2{,}15$ & $\le 2{,}79$ & $> 3{,}18$ \\
Mediocentros(MC) & $\le 2{,}10$ & $\le 2{,}60$ & $> 3{,}25$ \\
Delanteros  (DT) & $\le 2{,}25$ & $\le 2{,}75$ & $> 3{,}25$ \\
\hline
\end{tabular}
\caption{Umbrales de rendimiento por posición táctica.}
\label{tab:umbrales_rendimiento}
\end{table}

% ---------------------------------------------------------------------------
\subsubsection{Resultados obtenidos por demarcación}
\label{subsubsec:resultados_demarcacion}

Tras el entrenamiento y la optimización de hiperparámetros mediante
Grid Search, se presentan los resultados obtenidos para cada posición.
Se han evaluado modelos de regresión lineal regularizada (Ridge,
ElasticNet) y modelos basados en árboles (Random Forest, XGBoost).

\paragraph{A.\ Defensas (DF)}

El modelo Ridge destaca como el más preciso, situándose dentro del
rango de rendimiento aceptable con un $\mathrm{MAE} = 2{,}67$,
superando el objetivo establecido.

\begin{table}[H]
\centering
\small
\begin{tabular}{|l|r|r|r|p{3.2cm}|}
\hline
\textbf{Modelo} & \textbf{MAE} & \textbf{RMSE} & \textbf{Spear.} &
\textbf{Best Params} \\
\hline
Ridge         & 2{,}6748 & 3{,}3017 & 0{,}0889 & alpha: 1000{,}0 \\
ElasticNet    & 2{,}6756 & 3{,}3022 & 0{,}0874 & alpha: 0{,}05   \\
Random Forest & 2{,}6837 & 3{,}3098 & 0{,}0841 & max\_depth: 10  \\
XGBoost       & 2{,}7366 & 3{,}4026 & 0{,}0744 & lr: 0{,}1       \\
\hline
\end{tabular}
\caption{Resultados de los modelos evaluados para la posición Defensas (DF).}
\label{tab:resultados_df}
\end{table}

\paragraph{B.\ Delanteros (DT)}

En esta demarcación, ElasticNet obtuvo el mejor desempeño, quedando
muy cerca del umbral aceptable ($2{,}75$) con una correlación de
Spearman significativamente superior a la de los defensas
($\rho = 0{,}289$), lo que indica una mejor capacidad de \textit{ranking}.

\begin{table}[H]
\centering
\small
\begin{tabular}{|l|r|r|r|p{3.2cm}|}
\hline
\textbf{Modelo} & \textbf{MAE} & \textbf{RMSE} & \textbf{Spear.} &
\textbf{Best Params} \\
\hline
ElasticNet    & 2{,}7928 & 3{,}6878 & 0{,}2893 & alpha: 0{,}01  \\
Ridge         & 2{,}7981 & 3{,}6914 & 0{,}2859 & alpha: 100{,}0 \\
Random Forest & 2{,}8514 & 3{,}7055 & 0{,}2662 & max\_depth: 10 \\
XGBoost       & 2{,}9387 & 3{,}8463 & 0{,}2279 & n\_est: 300    \\
\hline
\end{tabular}
\caption{Resultados de los modelos evaluados para la posición Delanteros (DT).}
\label{tab:resultados_dt}
\end{table}

\paragraph{C.\ Mediocentros (MC)}

El modelo Ridge lidera la precisión en el centro del campo, mostrando
una estabilidad notable y situándose apenas $0{,}06$ puntos por encima
del objetivo «Aceptable». La correlación de Spearman ($\rho = 0{,}228$)
confirma la capacidad del modelo para ordenar correctamente a los
jugadores.

\begin{table}[H]
\centering
\small
\begin{tabular}{|l|r|r|r|p{3.2cm}|}
\hline
\textbf{Modelo} & \textbf{MAE} & \textbf{RMSE} & \textbf{Spear.} &
\textbf{Best Params} \\
\hline
Ridge         & 2{,}6650 & 3{,}4443 & 0{,}2277 & alpha: 250{,}0 \\
ElasticNet    & 2{,}6672 & 3{,}4457 & 0{,}2257 & alpha: 0{,}01  \\
Random Forest & 2{,}6701 & 3{,}4470 & 0{,}2259 & max\_depth: 10 \\
XGBoost       & 2{,}7621 & 3{,}5537 & 0{,}1566 & lr: 0{,}1      \\
\hline
\end{tabular}
\caption{Resultados de los modelos evaluados para la posición Mediocentros (MC).}
\label{tab:resultados_mc}
\end{table}

\paragraph{D.\ Porteros (PT)}

Para la posición de portero, ElasticNet ofrece el mejor ajuste,
cumpliendo prácticamente de forma exacta con el objetivo de rendimiento
aceptable establecido en $3{,}18$ ($\mathrm{MAE} = 3{,}19$).

\begin{table}[H]
\centering
\small
\begin{tabular}{|l|r|r|r|p{3.2cm}|}
\hline
\textbf{Modelo} & \textbf{MAE} & \textbf{RMSE} & \textbf{Spear.} &
\textbf{Best Params} \\
\hline
ElasticNet    & 3{,}1878 & 3{,}7598 & 0{,}1149 &  alpha: 0{,}05  \\
Ridge         & 3{,}2068 & 3{,}7750 & 0{,}0900 &  alpha: 250{,}0 \\
Random Forest & 3{,}2246 & 3{,}8273 & 0{,}0210 &  max\_depth: 20 \\
XGBoost       & 3{,}3537 & 4{,}0727 & -0{,}0407 & lr: 0{,}1      \\
\hline
\end{tabular}
\caption{Resultados de los modelos evaluados para la posición Porteros (PT).}
\label{tab:resultados_pt}
\end{table}

% ---------------------------------------------------------------------------
\subsubsection{Comparativa final y análisis del cumplimiento de objetivos}
\label{subsubsec:comparativa_final}

La tabla~\ref{tab:cumplimiento_objetivos} resume el cumplimiento de los
objetivos por posición utilizando el mejor modelo obtenido en cada caso.

\begin{table}[H]
\centering
\small
\begin{tabular}{|l|l|r|r|r|}
\hline
\textbf{Posición} & \textbf{Modelo} & \textbf{MAE} &
\textbf{Obj.} & \textbf{Estado} \\
\hline
Defensas     & Ridge      & 2{,}67 & 2{,}79 &  Cumplido   \\
Delanteros   & ElasticNet & 2{,}79 & 2{,}75 & $+0{,}04$              \\
Mediocentros & Ridge      & 2{,}66 & 2{,}60 & $+0{,}06$              \\
Porteros     & ElasticNet & 3{,}19 & 3{,}18 &  Cumplido   \\
\hline
\end{tabular}
\caption{Cumplimiento de los objetivos de rendimiento por posición.}
\label{tab:cumplimiento_objetivos}
\end{table}

\paragraph{Observaciones generales}

En líneas generales, los modelos de regresión lineal regularizada
(Ridge y ElasticNet) han demostrado una mayor capacidad de
generalización frente a modelos complejos como XGBoost en este
conjunto de datos. Los resultados son satisfactorios: en todas las
posiciones el $\mathrm{MAE}$ se mantiene significativamente alejado
del Umbral de Alerta (rango P75), validando la utilidad de los modelos
para la predicción de métricas en el contexto del estudio.

A pesar de que en las categorías de delanteros y mediocentros el
$\mathrm{MAE}$ se sitúa ligeramente por encima de la mediana fijada
como «objetivo aceptable», los modelos se consideran válidos y
funcionales para el propósito del estudio por las siguientes razones:

\paragraph{1.\ Capacidad de clasificación}

Las demarcaciones de delanteros ($\rho = 0{,}289$) y mediocentros
($\rho = 0{,}228$) presentan los coeficientes de Spearman más elevados
del proyecto. Esto indica que el modelo tiene una notable capacidad para
ordenar correctamente a los jugadores según su rendimiento real, lo cual
constituye el objetivo crítico de la herramienta. La capacidad de
\textit{ranking} resulta más relevante que la precisión absoluta cuando
el objetivo es identificar talento. Además, las desviaciones respecto al
objetivo son mínimas:

\begin{itemize}
  \item Delanteros: $\Delta\mathrm{MAE} = +0{,}04$
        (1{,}5\,\% por encima del objetivo).
  \item Mediocentros: $\Delta\mathrm{MAE} = +0{,}06$
        (2{,}3\,\% por encima del objetivo).
\end{itemize}

Ambas desviaciones se mantienen muy alejadas del Umbral de Alerta (P75),
con un margen de seguridad de $0{,}46$ y $0{,}59$ puntos respectivamente,
lo que demuestra la robustez de las predicciones.

\paragraph{2.\ Consistencia predictiva y generalización}

El hecho de que los modelos lineales regularizados (Ridge y ElasticNet)
sean los ganadores en todas las demarcaciones sugiere que el modelo ha
capturado tendencias generales del juego que no dependen del ruido ni
del \textit{overfitting}. Esta característica garantiza una predicción
real sobre datos no vistos, lo que es esencial para su aplicación en
competiciones o temporadas futuras.

Resulta altamente revelador que dichos modelos superen sistemáticamente
a algoritmos basados en árboles de decisión complejos (Random Forest y
XGBoost). Esto confirma que el rendimiento individual en el fútbol
contiene un componente intrínseco de aleatoriedad y ruido: mientras que
XGBoost tiende a sobreajustarse memorizando dicho ruido, los modelos
lineales demuestran una mayor capacidad de generalización al capturar
estrictamente la tendencia subyacente.

\paragraph{3.\ La paradoja varianza \textit{vs.}\ consistencia}

Los resultados exponen un patrón estadístico que refleja fielmente la
naturaleza estocástica del deporte:

\begin{itemize}
  \item \textbf{Atacantes (DT y MC):} presentan la mayor capacidad de
        clasificación ($\rho \approx 0{,}289$) pero un mayor error
        absoluto ($\mathrm{MAE}$). Es más sencillo ordenar la calidad
        de los atacantes porque los jugadores de élite destacan
        claramente en sus métricas base; sin embargo, predecir su
        puntuación exacta es complejo al depender de eventos de altísima
        varianza, como los goles o las asistencias.

  \item \textbf{Defensores y porteros (DF y PT):} cumplen con holgura
        los objetivos de $\mathrm{MAE}$ pero registran una capacidad de
        \textit{ranking} pobre ($\rho < 0{,}12$). Sus puntuaciones base
        son mucho más estables y predecibles (dependen del volumen de
        juego, los despejes o las porterías a cero), pero dada la alta
        homogeneidad en sus calificaciones empíricas, el modelo encuentra
        dificultades para establecer una jerarquía predictiva clara.
\end{itemize}

\paragraph{4.\ Alta penalización de características}

El análisis de los hiperparámetros óptimos arroja valores de penalización
notablemente altos, en especial en los modelos Ridge (p.\,ej.\
$\alpha = 1000{,}0$ en Defensas y $\alpha = 250{,}0$ en Mediocentros).
Esto refleja la necesidad del algoritmo de suprimir el impacto de
ciertas variables, indicando la presencia de multicolinealidad entre
los predictores extraídos mediante \textit{scraping}. Se confirma que
la regularización ha sido una técnica crítica y exitosa para filtrar
métricas redundantes que no aportaban valor real a la predicción final.

% ---------------------------------------------------------------------------
\subsubsection{Resumen de la validación final}
\label{subsubsec:resumen_validacion}

\begin{table}[H]
\centering
\small
\begin{tabular}{|l|l|c|c|p{4.2cm}|}
\hline
\textbf{Posición} & \textbf{Modelo} & \textbf{MAE} &
\textbf{Spearman} & \textbf{Evaluación} \\
\hline
Defensas     & Ridge      & 2{,}67 & 0{,}089 &
  Aceptable (error controlado)        \\
Delanteros   & ElasticNet & 2{,}79 & 0{,}289 &
  Excelente (buena capacidad ranking)  \\
Mediocentros & Ridge      & 2{,}66 & 0{,}228 &
  Bueno (error mínimo + ranking)      \\
Porteros     & ElasticNet & 3{,}19 & 0{,}115 &
  Aceptable (cumple exactamente)      \\
\hline
\end{tabular}
\caption{Resumen de la validación final por posición.}
\label{tab:resumen_validacion}
\end{table}

Se confirma la utilidad y robustez de la herramienta desarrollada.
Mientras que el $\mathrm{MAE}$ proporciona una medida de proximidad al
valor exacto predicho, el coeficiente de Spearman confirma que el modelo
comprende la jerarquía y estructura del rendimiento de los jugadores en
el mercado real de la Liga EA Sports, permitiendo:

\begin{enumerate}
  \item Identificar correctamente el talento infravalorado (jugadores
        con rendimiento superior al esperado según sus métricas base).
  \item Detectar jugadores en mala dinámica (con rendimiento inferior a
        sus métricas base).
  \item Proporcionar un \textit{ranking} fiable para procesos de
        fichaje y toma de decisiones.
\end{enumerate}

En consecuencia, los modelos desarrollados son adecuados para su
implementación en sistemas de recomendación de transferencias y en el
análisis predictivo del rendimiento en el contexto de la competición
analizada.